% ── ABSTRACT (~300 words) ─────────────────────────────────────────────────────

Three independent intellectual traditions --- modern quantum physics, the
process philosophy of Alfred North Whitehead, and the Bahá'í writings of
Bahá'u'lláh and 'Abdu'l-Bahá --- converge on a structurally identical
description of reality: that existence is constituted by relationship, that
properties emerge only through relational interaction, and that the force
sustaining those relationships is, in each tradition's vocabulary, love.

This paper makes three related arguments.
First, substance ontology is \emph{formally incompatible} with quantum mechanics:
Bell's theorem \citep{Bell1964} eliminates intrinsic properties, making
Whitehead's process ontology not a permitted interpretation but the only
internally consistent metaphysical framework available.
Second, the Bahá'í writings constitute an \emph{independent prior instantiation}
of process ontology, articulating its essential claims forty years before
\citet{Whitehead1929} named the framework philosophically and a century before
physics confirmed it experimentally.
Third, these two findings together establish a \emph{triadic convergence}:
three traditions built by three different civilizations, using three different
methodologies, across 170 years, detect the same feature of reality.

The convergence is formalized through six structural isomorphisms (relational
ontology, nonlocal interconnection, observer-dependence, decoherence, holographic
order, and motion as life) and the Relational Coherence Theorem (\RCT),
developed in a companion physics paper, which translates the philosophical claim
into testable quantum-mechanical predictions.
The argument employs an overlapping consensus architecture: each tradition
validates its own thesis independently; the convergence is not a merged view
but an independent confirmation by separate instruments of the same underlying
structure of reality.
